\chapter{Chapter09 — One Bit Hypercontractivity}
%%% generated by the blueprint pipeline from FABL.Chapter09.OneBitHypercontractivity
\section{Overview}

Declarations of the Lean module \texttt{FABL.Chapter09.OneBitHypercontractivity}.

\section{Declarations}

\begin{definition}[Uniform maximum norm]
\lean{FABL.uniformLInfinityNorm}
\label{FABL.uniformLInfinityNorm}
\leanok
Let $\Omega$ be a finite nonempty set, and let $f:\Omega\to\mathbb R$. Define the
uniform $L^\infty$ norm of $f$ to be $\max_{x\in\Omega}|f(x)|$.
\end{definition}

\begin{definition}[Uniform norm with an extended exponent]
\lean{FABL.uniformLpNormENNReal}
\label{FABL.uniformLpNormENNReal}
\leanok
\uses{FABL.uniformLInfinityNorm, FABL.uniformLpNorm}
Let $\Omega$ be a finite nonempty set, let $q\in[0,\infty]$, and let $f:\Omega\to\mathbb
R$. Define its uniform norm with exponent $q$ to be its uniform $L^\infty$ norm when
$q=\infty$, and otherwise its uniform $L^q$ norm with exponent equal to the real value
of $q$.
\end{definition}

\begin{theorem}[Finite-exponent norm under coercion]
\lean{FABL.uniformLpNormENNReal_coe}
\label{FABL.uniformLpNormENNReal_coe}
\leanok
\uses{FABL.uniformLpNorm, FABL.uniformLpNormENNReal}
\difficulty{1}
For a finite set $\Omega$, a nonnegative real number $q$, and a function
$f:\Omega\to\mathbb R$, the extended-exponent uniform norm at the coercion of $q$ to the
extended nonnegative reals equals the uniform $q$-norm.
\end{theorem}

\begin{theorem}[Infinite-exponent uniform norm]
\lean{FABL.uniformLpNormENNReal_top}
\label{FABL.uniformLpNormENNReal_top}
\leanok
\uses{FABL.uniformLInfinityNorm, FABL.uniformLpNormENNReal}
\difficulty{1}
For a finite nonempty set $\Omega$ and a function $f:\Omega\to\mathbb R$, the
extended-exponent uniform norm of $f$ at exponent $\infty$ equals its supremum norm.
\end{theorem}

\begin{theorem}[Nonnegativity of the supremum norm]
\lean{FABL.uniformLInfinityNorm_nonneg}
\label{FABL.uniformLInfinityNorm_nonneg}
\leanok
\uses{FABL.uniformLInfinityNorm}
\difficulty{3}
For a finite nonempty set $\Omega$ and a function $f:\Omega\to\mathbb R$, the supremum
norm of $f$ is nonnegative.
\end{theorem}

\begin{theorem}[Pointwise bound by the uniform norm]
\lean{FABL.abs_le_uniformLInfinityNorm}
\label{FABL.abs_le_uniformLInfinityNorm}
\leanok
\uses{FABL.uniformLInfinityNorm}
\difficulty{3}
Let $\Omega$ be a finite nonempty set, let $f\colon\Omega\to\mathbb R$, and let
$x\in\Omega$. Then $|f(x)|\leq \max_{y\in\Omega}|f(y)|$.
\end{theorem}

\begin{theorem}[Supremum norm on the Boolean domain]
\lean{FABL.uniformLInfinityNorm_bool}
\label{FABL.uniformLInfinityNorm_bool}
\leanok
\uses{FABL.uniformLInfinityNorm}
\difficulty{3}
For a function $f:\{\mathrm{false},\mathrm{true}\}\to\mathbb R$, its supremum norm is
$\max\{|f(\mathrm{false})|,|f(\mathrm{true})|\}$.
\end{theorem}

\begin{theorem}[Nonnegativity of the extended-exponent uniform norm]
\lean{FABL.uniformLpNormENNReal_nonneg}
\label{FABL.uniformLpNormENNReal_nonneg}
\leanok
\uses{FABL.uniformLInfinityNorm_nonneg, FABL.uniformLpNormENNReal, FABL.uniformLpNorm_nonneg}
\difficulty{4}
For a finite nonempty set $\Omega$, an extended nonnegative real exponent $q$, and a
function $f:\Omega\to\mathbb R$, the extended-exponent uniform norm of $f$ is
nonnegative.
\end{theorem}

\begin{theorem}[Scaling the uniform norm]
\lean{FABL.uniformLpNorm_const_mul}
\label{FABL.uniformLpNorm_const_mul}
\leanok
\uses{FABL.uniformLpNorm}
\difficulty{4}
For a finite set $\Omega$, real numbers $p,c$ with $p>0$, and a function
$f:\Omega\to\mathbb{R}$, the uniform $L^p$ quantity of $cf$ equals $|c|$ times the
uniform $L^p$ quantity of $f$.
\end{theorem}

\begin{theorem}[Scaling the supremum norm]
\lean{FABL.uniformLInfinityNorm_const_mul}
\label{FABL.uniformLInfinityNorm_const_mul}
\leanok
\uses{FABL.uniformLInfinityNorm}
\difficulty{3}
For a finite nonempty set $\Omega$, a real number $c$, and a function
$f:\Omega\to\mathbb R$, the supremum norm of $x\mapsto c f(x)$ equals $|c|$ times the
supremum norm of $f$.
\end{theorem}

\begin{theorem}[Scaling the extended-exponent uniform norm]
\lean{FABL.uniformLpNormENNReal_const_mul}
\label{FABL.uniformLpNormENNReal_const_mul}
\leanok
\uses{FABL.uniformLInfinityNorm_const_mul, FABL.uniformLpNormENNReal, FABL.uniformLpNorm_const_mul}
\difficulty{4}
For a finite nonempty set $\Omega$, an extended nonnegative real exponent $q\ne 0$, a
real number $c$, and a function $f:\Omega\to\mathbb R$, the extended-exponent uniform
norm of $x\mapsto c f(x)$ equals $|c|$ times that of $f$.
\end{theorem}

\begin{definition}[Hypercontractive random variable]
\lean{FABL.IsHypercontractiveRandomVariable}
\statementsource{Analysis-of-Boolean-Functions-by-Ryan-ODonnell}{
  pdf-ffbe9c7b3306-p256-b003,
  pdf-ffbe9c7b3306-p259-b004,
  pdf-ffbe9c7b3306-p259-b006
}
\label{FABL.IsHypercontractiveRandomVariable}
\leanok
\uses{FABL.uniformLpNormENNReal}
Let $\Omega$ be a finite nonempty probability space with the uniform distribution, let
$p,q\in[0,\infty]$, let $\rho\in\mathbb{R}$, and let $X:\Omega\to\mathbb{R}$. The random
variable $X$ is hypercontractive with parameters $(p,q,\rho)$ if, for every
$a,b\in\mathbb{R}$, the uniform $L^q$ norm of $a+\rho bX$ is at most the uniform $L^p$
norm of $a+bX$.
\end{definition}

\begin{theorem}[Scaling a hypercontractive variable]
\lean{FABL.IsHypercontractiveRandomVariable.const_mul}
\label{FABL.IsHypercontractiveRandomVariable.const_mul}
\leanok
\uses{FABL.IsHypercontractiveRandomVariable}
\difficulty{1}
Let $\\Omega$ be a finite nonempty set, let $X:\\Omega\\to\\mathbb R$, let
$p,q\\in[0,\\infty]$, and let $\\rho,c\\in\\mathbb R$. If for every $a,b\\in\\mathbb R$
the uniform $L^q$ norm of $a+\\rho bX$ is at most the uniform $L^p$ norm of $a+bX$, then
the same inequality holds with $X$ replaced by the function $x\\mapsto cX(x)$. At
exponent $\\infty$, the uniform norm means the maximum absolute value on $\\Omega$.
\end{theorem}

\begin{theorem}[Hypercontractivity from normalized tests]
\lean{FABL.IsHypercontractiveRandomVariable.of_normalized}
\label{FABL.IsHypercontractiveRandomVariable.of_normalized}
\leanok
\uses{FABL.IsHypercontractiveRandomVariable, FABL.uniformLpNormENNReal, FABL.uniformLpNormENNReal_const_mul}
\difficulty{5}
Let $\Omega$ be a finite nonempty set, let $p,q\in[0,\infty]$ be nonzero, let
$\rho\in\mathbb R$, and let $X:\Omega\to\mathbb R$. Suppose that for every
$a,b\in\mathbb R$ with $a=1$ or $b=1$, the uniform $q$-norm of $x\mapsto a+\rho bX(x)$
is at most the uniform $p$-norm of $x\mapsto a+bX(x)$. Then this inequality holds for
every $a,b\in\mathbb R$.
\end{theorem}

\begin{theorem}[Contraction parameter times norm bound]
\lean{FABL.IsHypercontractiveRandomVariable.rho_mul_norm_le}
\label{FABL.IsHypercontractiveRandomVariable.rho_mul_norm_le}
\leanok
\uses{FABL.IsHypercontractiveRandomVariable, FABL.uniformLpNormENNReal, FABL.uniformLpNormENNReal_const_mul}
\difficulty{3}
Let $\Omega$ be a finite nonempty set, let $p,q\in[0,\infty]$ with $q\ne0$, let
$\rho\ge0$, and let $X:\Omega\to\mathbb R$. Suppose the uniform $q$-norm of $x\mapsto
a+\rho bX(x)$ is at most the uniform $p$-norm of $x\mapsto a+bX(x)$ for every pair of
real numbers $a,b$. Then $\rho$ times the uniform $q$-norm of $X$ is at most its uniform
$p$-norm.
\end{theorem}

\begin{theorem}[Contraction parameter bounded by a norm ratio]
\lean{FABL.IsHypercontractiveRandomVariable.rho_le_norm_ratio}
\proofsource{Analysis-of-Boolean-Functions-by-Ryan-ODonnell}{pdf-ffbe9c7b3306-p272-b001}
\label{FABL.IsHypercontractiveRandomVariable.rho_le_norm_ratio}
\leanok
\uses{FABL.IsHypercontractiveRandomVariable, FABL.IsHypercontractiveRandomVariable.rho_mul_norm_le, FABL.uniformLpNormENNReal}
\difficulty{1}
Let $\Omega$ be a finite nonempty set, let $p,q\in[0,\infty]$ with $q\ne0$, let
$\rho\ge0$, and let $X:\Omega\to\mathbb R$. Suppose the uniform $q$-norm of $x\mapsto
a+\rho bX(x)$ is at most the uniform $p$-norm of $x\mapsto a+bX(x)$ for all real $a,b$,
and suppose the uniform $q$-norm of $X$ is positive. Then $\rho$ is at most the uniform
$p$-norm of $X$ divided by its uniform $q$-norm.
\end{theorem}

\begin{theorem}[Norm of an affine sign function]
\lean{FABL.uniformLpNorm_two_affineSign}
\label{FABL.uniformLpNorm_two_affineSign}
\leanok
\uses{FABL.Sign, FABL.signValue, FABL.uniformLpNorm, FABL.uniformLpNorm_two_eq_sqrt_expect_sq}
\difficulty{3}
For all real numbers $a$ and $b$, the uniform $2$-norm of the function on $\{-1,1\}$
given by $s\mapsto a+bs$ is $\sqrt{a^2+b^2}$.
\end{theorem}

\begin{theorem}[Affine sign supremum norm]
\lean{FABL.uniformLInfinityNorm_affineSign}
\label{FABL.uniformLInfinityNorm_affineSign}
\leanok
\uses{FABL.Sign, FABL.signValue, FABL.uniformLInfinityNorm}
\difficulty{4}
For real numbers $a$ and $b$, the supremum of $|a+b s|$ over the two integer signs
$s\in\{-1,1\}$ is $\max\{|a+b|,|a-b|\}$.
\end{theorem}

\begin{theorem}[Uniform norm of an affine sign function]
\lean{FABL.uniformLpNorm_affineSign}
\label{FABL.uniformLpNorm_affineSign}
\leanok
\uses{FABL.Sign, FABL.signValue, FABL.uniformLpNorm}
\difficulty{3}
For real numbers $p,a,b$, let a sign be either $-1$ or $1$, with its usual real value.
The uniform $L^p$ quantity of the function sending a sign $s$ to $a+bs$ is
$\left((|a+b|^p+|a-b|^p)/2\right)^{1/p}$.
\end{theorem}

\begin{theorem}[Sixth moment of the uniform norm]
\lean{FABL.uniformLpNorm_six_pow_eq_expect_six}
\label{FABL.uniformLpNorm_six_pow_eq_expect_six}
\leanok
\uses{FABL.uniformLpNorm}
\difficulty{5}
Let $\Omega$ be a finite set and let $f\colon\Omega\to\mathbb R$. The sixth power of the
uniform $L^6$ norm of $f$ equals the uniform average of $f^6$.
\end{theorem}

\begin{theorem}[Fourth moment of an affine random sign]
\lean{FABL.expect_affineSign_fourth}
\label{FABL.expect_affineSign_fourth}
\leanok
\uses{FABL.Sign, FABL.signValue}
\difficulty{3}
Let $a,b\in\mathbb{R}$ and let $S$ be uniformly distributed on $\{-1,1\}$. Then
$\mathbb{E}(a+bS)^4=a^4+6a^2b^2+b^4$.
\end{theorem}

\begin{theorem}[Sixth moment of an affine random sign]
\lean{FABL.expect_affineSign_sixth}
\proofsource{Analysis-of-Boolean-Functions-by-Ryan-ODonnell}{pdf-ffbe9c7b3306-p257-b004}
\label{FABL.expect_affineSign_sixth}
\leanok
\uses{FABL.Sign, FABL.signValue}
\difficulty{3}
Let $a,b\in\mathbb{R}$ and let $S$ be uniformly distributed on $\{-1,1\}$. Then
$\mathbb{E}(a+bS)^6=a^6+15a^4b^2+15a^2b^4+b^6$.
\end{theorem}

\begin{theorem}[Average rpow lower quadratic]
\lean{FABL.average_rpow_lower_quadratic}
\statementsource{Analysis-of-Boolean-Functions-by-Ryan-ODonnell}{pdf-ffbe9c7b3306-p258-b003}
\label{FABL.average_rpow_lower_quadratic}
\leanok
\difficulty{7}
% dataset-pipeline: stub -- prose pending build_informal_statements.py
Statement pending regeneration.
\end{theorem}

\begin{theorem}[Normalized p two affine]
\lean{FABL.normalized_pTwo_affine}
\statementsource{Analysis-of-Boolean-Functions-by-Ryan-ODonnell}{pdf-ffbe9c7b3306-p258-b003}
\label{FABL.normalized_pTwo_affine}
\leanok
\uses{FABL.average_rpow_lower_quadratic}
\difficulty{4}
% dataset-pipeline: stub -- prose pending build_informal_statements.py
Statement pending regeneration.
\end{theorem}

\begin{theorem}[Affine p two of abs le]
\lean{FABL.affine_pTwo_of_abs_le}
\statementsource{Analysis-of-Boolean-Functions-by-Ryan-ODonnell}{pdf-ffbe9c7b3306-p258-b003}
\label{FABL.affine_pTwo_of_abs_le}
\leanok
\uses{FABL.normalized_pTwo_affine}
\difficulty{5}
% dataset-pipeline: stub -- prose pending build_informal_statements.py
Statement pending regeneration.
\end{theorem}

\begin{theorem}[Affine form of Hanner’s inequality]
\lean{FABL.affine_pTwo_hanner}
\proofsource{Analysis-of-Boolean-Functions-by-Ryan-ODonnell}{pdf-ffbe9c7b3306-p258-b003}
\label{FABL.affine_pTwo_hanner}
\leanok
\uses{FABL.affine_pTwo_of_abs_le}
\difficulty{4}
For real numbers $p,a,b$ satisfying $1\le p\le 2$, one has
\[
\sqrt{a^2+(p-1)b^2}\le\left(\frac{|a+b|^p+|a-b|^p}{2}\right)^{1/p}.
\]
\end{theorem}

\begin{theorem}[One bit p two]
\lean{FABL.oneBit_pTwo}
\statementsource{Analysis-of-Boolean-Functions-by-Ryan-ODonnell}{
  pdf-ffbe9c7b3306-p248-b001,
  pdf-ffbe9c7b3306-p258-b003
}
\label{FABL.oneBit_pTwo}
\leanok
\uses{FABL.Sign, FABL.affine_pTwo_hanner, FABL.signValue, FABL.uniformLpNorm, FABL.uniformLpNorm_affineSign, FABL.uniformLpNorm_two_affineSign}
\difficulty{3}
% dataset-pipeline: stub -- prose pending build_informal_statements.py
Statement pending regeneration.
\end{theorem}

\begin{theorem}[One bit p two sharp]
\lean{FABL.oneBit_pTwo_sharp}
\statementsource{Analysis-of-Boolean-Functions-by-Ryan-ODonnell}{pdf-ffbe9c7b3306-p258-b003}
\label{FABL.oneBit_pTwo_sharp}
\leanok
\uses{FABL.Sign, FABL.oneBit_pTwo, FABL.signValue, FABL.uniformLpNorm}
\difficulty{1}
% dataset-pipeline: stub -- prose pending build_informal_statements.py
Statement pending regeneration.
\end{theorem}

\begin{theorem}[Uniform sign p two hypercontractive]
\lean{FABL.uniformSign_pTwoHypercontractive}
\statementsource{Analysis-of-Boolean-Functions-by-Ryan-ODonnell}{
  pdf-ffbe9c7b3306-p248-b001,
  pdf-ffbe9c7b3306-p259-b006
}
\label{FABL.uniformSign_pTwoHypercontractive}
\leanok
\uses{FABL.IsHypercontractiveRandomVariable, FABL.oneBit_pTwo_sharp, FABL.signValue, FABL.uniformLpNormENNReal, FABL.uniformLpNormENNReal_coe}
\difficulty{3}
% dataset-pipeline: stub -- prose pending build_informal_statements.py
Statement pending regeneration.
\end{theorem}

\begin{theorem}[One bit two four]
\lean{FABL.oneBit_twoFour}
\statementsource{Analysis-of-Boolean-Functions-by-Ryan-ODonnell}{pdf-ffbe9c7b3306-p253-b002}
\label{FABL.oneBit_twoFour}
\leanok
\uses{FABL.Sign, FABL.expect_affineSign_fourth, FABL.signValue, FABL.uniformLpNorm, FABL.uniformLpNorm_four_pow_eq_expect_fourth, FABL.uniformLpNorm_nonneg, FABL.uniformLpNorm_two_affineSign}
\difficulty{5}
% dataset-pipeline: stub -- prose pending build_informal_statements.py
Statement pending regeneration.
\end{theorem}

\begin{definition}[Sixth noise root]
\lean{FABL.twoSixNoiseRoot}
\statementsource{Analysis-of-Boolean-Functions-by-Ryan-ODonnell}{pdf-ffbe9c7b3306-p257-b003}
\label{FABL.twoSixNoiseRoot}
\leanok
Define the six-noise root to be $\sqrt{1/5}$.
\end{definition}

\begin{theorem}[Nonnegativity of the noise root]
\lean{FABL.twoSixNoiseRoot_nonneg}
\label{FABL.twoSixNoiseRoot_nonneg}
\leanok
\uses{FABL.twoSixNoiseRoot}
\difficulty{1}
The nonnegative square root of $1/5$ is nonnegative.
\end{theorem}

\begin{theorem}[Square of the noise root]
\lean{FABL.twoSixNoiseRoot_sq}
\label{FABL.twoSixNoiseRoot_sq}
\leanok
\uses{FABL.twoSixNoiseRoot}
\difficulty{1}
The square of the nonnegative square root of $1/5$ equals $1/5$.
\end{theorem}

\begin{theorem}[Sixth-moment norm comparison]
\lean{FABL.oneBit_twoSix}
\proofsource{Analysis-of-Boolean-Functions-by-Ryan-ODonnell}{
  pdf-ffbe9c7b3306-p257-b003,
  pdf-ffbe9c7b3306-p257-b004
}
\label{FABL.oneBit_twoSix}
\leanok
\uses{FABL.Sign, FABL.expect_affineSign_sixth, FABL.signValue, FABL.uniformLpNorm, FABL.uniformLpNorm_nonneg, FABL.uniformLpNorm_six_pow_eq_expect_six, FABL.uniformLpNorm_two_affineSign}
\difficulty{5}
Let $a,b,\rho\in\mathbb R$ satisfy $\rho^2\le 1/5$. For a uniformly chosen sign
$\varepsilon\in\{-1,1\}$, the sixth-norm of $a+\rho b\varepsilon$ is at most the
$2$-norm of $a+b\varepsilon$.
\end{theorem}

\begin{theorem}[Sharp sixth-norm comparison]
\lean{FABL.oneBit_twoSix_sharp}
\proofsource{Analysis-of-Boolean-Functions-by-Ryan-ODonnell}{
  pdf-ffbe9c7b3306-p257-b003,
  pdf-ffbe9c7b3306-p257-b004
}
\label{FABL.oneBit_twoSix_sharp}
\leanok
\uses{FABL.Sign, FABL.oneBit_twoSix, FABL.signValue, FABL.twoSixNoiseRoot, FABL.twoSixNoiseRoot_sq, FABL.uniformLpNorm}
\difficulty{1}
Let $a,b\in\mathbb R$. For a uniformly chosen sign $\varepsilon\in\{-1,1\}$, the
sixth-norm of $a+\sqrt{1/5}\,b\varepsilon$ is at most the $2$-norm of $a+b\varepsilon$.
\end{theorem}

\begin{theorem}[Necessary sixth-norm noise bound]
\lean{FABL.oneBit_twoSix_necessary_sq}
\proofsource{Analysis-of-Boolean-Functions-by-Ryan-ODonnell}{pdf-ffbe9c7b3306-p257-b004}
\label{FABL.oneBit_twoSix_necessary_sq}
\leanok
\uses{FABL.Sign, FABL.expect_affineSign_sixth, FABL.signValue, FABL.uniformLpNorm, FABL.uniformLpNorm_nonneg, FABL.uniformLpNorm_six_pow_eq_expect_six, FABL.uniformLpNorm_two_affineSign}
\difficulty{6}
Let $\rho\in\mathbb R$. If, for every $a,b\in\mathbb R$, the sixth-norm of $a+\rho
b\varepsilon$ is at most the $2$-norm of $a+b\varepsilon$ for a uniformly chosen sign
$\varepsilon\in\{-1,1\}$, then $\rho^2\le 1/5$.
\end{theorem}

\begin{theorem}[Exact sixth-norm noise threshold]
\lean{FABL.oneBit_twoSix_iff}
\proofsource{Analysis-of-Boolean-Functions-by-Ryan-ODonnell}{
  pdf-ffbe9c7b3306-p257-b003,
  pdf-ffbe9c7b3306-p257-b004
}
\label{FABL.oneBit_twoSix_iff}
\leanok
\uses{FABL.Sign, FABL.oneBit_twoSix, FABL.oneBit_twoSix_necessary_sq, FABL.signValue, FABL.twoSixNoiseRoot, FABL.twoSixNoiseRoot_nonneg, FABL.twoSixNoiseRoot_sq, FABL.uniformLpNorm}
\difficulty{4}
Let $\rho\in\mathbb R$ satisfy $0\le\rho$. The sixth-norm of $a+\rho b\varepsilon$ is at
most the $2$-norm of $a+b\varepsilon$ for every $a,b\in\mathbb R$ and a uniformly chosen
sign $\varepsilon\in\{-1,1\}$ if and only if $\rho\le\sqrt{1/5}$.
\end{theorem}

\begin{theorem}[Two to four hypercontractivity]
\lean{FABL.uniformSign_twoFourHypercontractive}
\proofsource{Analysis-of-Boolean-Functions-by-Ryan-ODonnell}{
  pdf-ffbe9c7b3306-p259-b006,
  pdf-ffbe9c7b3306-p253-b002
}
\label{FABL.uniformSign_twoFourHypercontractive}
\leanok
\uses{FABL.IsHypercontractiveRandomVariable, FABL.oneBit_twoFour, FABL.signValue, FABL.twoFourNoiseRoot, FABL.twoFourNoiseRoot_sq, FABL.uniformLpNormENNReal}
\difficulty{3}
Let $S=\{-1,1\}$ have the uniform probability distribution, and let $X(s)=s$. For every
$a,b\in\mathbb{R}$, the uniform $L^4$ norm of $a+bX/\sqrt{3}$ is at most the uniform
$L^2$ norm of $a+bX$.
\end{theorem}

\begin{theorem}[Two to six hypercontractivity]
\lean{FABL.uniformSign_twoSixHypercontractive}
\proofsource{Analysis-of-Boolean-Functions-by-Ryan-ODonnell}{
  pdf-ffbe9c7b3306-p257-b003,
  pdf-ffbe9c7b3306-p257-b004
}
\label{FABL.uniformSign_twoSixHypercontractive}
\leanok
\uses{FABL.IsHypercontractiveRandomVariable, FABL.oneBit_twoSix_sharp, FABL.signValue, FABL.twoSixNoiseRoot, FABL.uniformLpNormENNReal}
\difficulty{3}
Let $S=\{-1,1\}$ have the uniform probability distribution, and let $X(s)=s$. For every
$a,b\in\mathbb{R}$, the uniform $L^6$ norm of $a+bX/\sqrt{5}$ is at most the uniform
$L^2$ norm of $a+bX$.
\end{theorem}

\begin{theorem}[One bit four third two]
\lean{FABL.oneBit_fourThirdTwo}
\statementsource{Analysis-of-Boolean-Functions-by-Ryan-ODonnell}{
  pdf-ffbe9c7b3306-p254-b002,
  pdf-ffbe9c7b3306-p248-b001,
  pdf-ffbe9c7b3306-p259-b006
}
\label{FABL.oneBit_fourThirdTwo}
\leanok
\uses{FABL.Sign, FABL.oneBit_twoFour, FABL.signValue, FABL.twoFourNoiseRoot, FABL.twoFourNoiseRoot_sq, FABL.uniformInner, FABL.uniformInner_le_uniformLpNorm_fourThird_mul_four, FABL.uniformLpNorm, FABL.uniformLpNorm_nonneg, FABL.uniformLpNorm_two_sq_eq_expect_sq}
\difficulty{5}
% dataset-pipeline: stub -- prose pending build_informal_statements.py
Statement pending regeneration.
\end{theorem}

\begin{theorem}[Hypercontractivity of a uniform sign]
\lean{FABL.uniformSign_fourThirdTwoHypercontractive}
\proofsource{Analysis-of-Boolean-Functions-by-Ryan-ODonnell}{
  pdf-ffbe9c7b3306-p254-b002,
  pdf-ffbe9c7b3306-p254-b003
}
\label{FABL.uniformSign_fourThirdTwoHypercontractive}
\leanok
\uses{FABL.IsHypercontractiveRandomVariable, FABL.oneBit_fourThirdTwo, FABL.signValue, FABL.twoFourNoiseRoot, FABL.uniformLpNormENNReal, FABL.uniformLpNormENNReal_coe}
\difficulty{3}
Let $S$ be the two-element set of signs $\{-1,1\}$, with the uniform probability
measure, and let $X(s)=s$. For every pair of real numbers $a,b$, the uniform $L^2$ norm
of $a+bX/\sqrt{3}$ is at most the uniform $L^{4/3}$ norm of $a+bX$.
\end{theorem}

\begin{theorem}[One-bit supremum norm bound]
\lean{FABL.oneBit_twoTop}
\proofsource{Analysis-of-Boolean-Functions-by-Ryan-ODonnell}{
  pdf-ffbe9c7b3306-p247-b006,
  pdf-ffbe9c7b3306-p259-b006
}
\label{FABL.oneBit_twoTop}
\leanok
\uses{FABL.Sign, FABL.signValue, FABL.uniformLInfinityNorm, FABL.uniformLInfinityNorm_affineSign, FABL.uniformLpNorm, FABL.uniformLpNorm_two_affineSign}
\difficulty{4}
Let $a,b\in\mathbb R$. The supremum of $|a+0\cdot b\varepsilon|$ over signs
$\varepsilon\in\{-1,1\}$ is at most the $2$-norm of $a+b\varepsilon$ for a uniformly
chosen sign.
\end{theorem}

\begin{theorem}[Two to infinity norm bound]
\lean{FABL.uniformSign_twoTopHypercontractive}
\proofsource{Analysis-of-Boolean-Functions-by-Ryan-ODonnell}{
  pdf-ffbe9c7b3306-p259-b004,
  pdf-ffbe9c7b3306-p259-b006,
  pdf-ffbe9c7b3306-p247-b006
}
\label{FABL.uniformSign_twoTopHypercontractive}
\leanok
\uses{FABL.IsHypercontractiveRandomVariable, FABL.oneBit_twoTop, FABL.signValue, FABL.uniformLpNormENNReal, FABL.uniformLpNormENNReal_top}
\difficulty{3}
Let $S=\{-1,1\}$ have the uniform probability distribution, and let $X(s)=s$. For every
$a,b\in\mathbb{R}$, the uniform supremum norm of the constant function with value $a$ is
at most the uniform $L^2$ norm of $a+bX$.
\end{theorem}

\begin{theorem}[Hölder inequality for the uniform inner product]
\lean{FABL.uniformInner_le_uniformLpNorm_mul}
\label{FABL.uniformInner_le_uniformLpNorm_mul}
\leanok
\uses{FABL.uniformInner, FABL.uniformLpNorm}
\difficulty{5}
Let $\Omega$ be a finite nonempty set, let $p,q\in\mathbb R$ be Hölder conjugates, and
let $f,g:\Omega\to\mathbb R$. Then the uniform inner product satisfies $\langle
f,g\rangle_u\leq \|f\|_{p,u}\,\|g\|_{q,u}$, where the uniform $r$-norm is the uniform
expectation of $|h|^r$ raised to the power $1/r$.
\end{theorem}

\begin{theorem}[Expectation bounded by a uniform $L^q$ norm]
\lean{FABL.abs_expect_le_uniformLpNorm}
\label{FABL.abs_expect_le_uniformLpNorm}
\leanok
\uses{FABL.uniformLpNorm}
\difficulty{4}
Let $\Omega$ be a finite nonempty set, let $q\geq1$, and let
$f\colon\Omega\to\mathbb{R}$. The absolute value of the uniform average of $f$ on
$\Omega$ is at most its uniform $L^q$ norm, defined as
$\bigl(\mathbb{E}_{x\in\Omega}|f(x)|^q\bigr)^{1/q}$.
\end{theorem}

\begin{theorem}[Triangle inequality for the uniform norm]
\lean{FABL.uniformLpNorm_add_le}
\label{FABL.uniformLpNorm_add_le}
\leanok
\uses{FABL.uniformLpNorm}
\difficulty{5}
For a finite nonempty set $\Omega$, a real number $q\ge 1$, and functions
$f,g:\Omega\to\mathbb{R}$, the uniform $L^q$ quantity of $f+g$ is at most the sum of the
uniform $L^q$ quantities of $f$ and $g$.
\end{theorem}

\begin{theorem}[Mean zero qq hypercontractive real]
\lean{FABL.meanZero_qq_hypercontractive_real}
\statementsource{Analysis-of-Boolean-Functions-by-Ryan-ODonnell}{pdf-ffbe9c7b3306-p272-b003}
\label{FABL.meanZero_qq_hypercontractive_real}
\leanok
\uses{FABL.abs_expect_le_uniformLpNorm, FABL.uniformLpNorm, FABL.uniformLpNorm_add_le, FABL.uniformLpNorm_const_mul}
\difficulty{5}
% dataset-pipeline: stub -- prose pending build_informal_statements.py
Statement pending regeneration.
\end{theorem}

\begin{theorem}[Mean zero qq hypercontractive]
\lean{FABL.meanZero_qq_hypercontractive}
\statementsource{Analysis-of-Boolean-Functions-by-Ryan-ODonnell}{pdf-ffbe9c7b3306-p272-b003}
\label{FABL.meanZero_qq_hypercontractive}
\leanok
\uses{FABL.IsHypercontractiveRandomVariable, FABL.meanZero_qq_hypercontractive_real, FABL.uniformLpNormENNReal_coe}
\difficulty{1}
% dataset-pipeline: stub -- prose pending build_informal_statements.py
Statement pending regeneration.
\end{theorem}

\begin{theorem}[Composition of hypercontractive inequalities]
\lean{FABL.IsHypercontractiveRandomVariable.comp}
\label{FABL.IsHypercontractiveRandomVariable.comp}
\leanok
\uses{FABL.IsHypercontractiveRandomVariable, FABL.uniformLpNormENNReal}
\difficulty{3}
Let $\\Omega$ be a finite nonempty set, let $X:\\Omega\\to\\mathbb R$, let
$p,r,q\\in[0,\\infty]$, and let $\\sigma,\\rho\\in\\mathbb R$. Suppose that for every
$a,b\\in\\mathbb R$, the uniform $L^r$ norm of $a+\\sigma bX$ is at most the uniform
$L^p$ norm of $a+bX$, and the uniform $L^q$ norm of $a+\\rho bX$ is at most the uniform
$L^r$ norm of $a+bX$. Then for every $a,b\\in\\mathbb R$, the uniform $L^q$ norm of
$a+\\rho\\sigma bX$ is at most the uniform $L^p$ norm of $a+bX$. Here the uniform norm
at exponent $\\infty$ is the maximum absolute value on $\\Omega$.
\end{theorem}

\begin{theorem}[Mono rho of mean zero]
\lean{FABL.IsHypercontractiveRandomVariable.mono_rho_of_meanZero}
\statementsource{Analysis-of-Boolean-Functions-by-Ryan-ODonnell}{
  pdf-ffbe9c7b3306-p256-b004,
  pdf-ffbe9c7b3306-p272-b003
}
\label{FABL.IsHypercontractiveRandomVariable.mono_rho_of_meanZero}
\leanok
\uses{FABL.IsHypercontractiveRandomVariable, FABL.IsHypercontractiveRandomVariable.comp, FABL.meanZero_qq_hypercontractive}
\difficulty{3}
% dataset-pipeline: stub -- prose pending build_informal_statements.py
Statement pending regeneration.
\end{theorem}

\begin{theorem}[A dual function attaining the conjugate norm]
\lean{FABL.exists_uniformLpNorm_dual}
\label{FABL.exists_uniformLpNorm_dual}
\leanok
\uses{FABL.uniformInner, FABL.uniformLpNorm}
\difficulty{7}
Let $\Omega$ be a finite nonempty set, let $p,q\in\mathbb R$ be Hölder conjugates, and
let $f:\Omega\to\mathbb R$. There is a function $g:\Omega\to\mathbb R$ whose uniform
$p$-norm is at most $1$ and whose uniform inner product with $f$ equals the uniform
$q$-norm of $f$.
\end{theorem}

\begin{theorem}[Dual bound for a self-adjoint operator]
\lean{FABL.selfAdjoint_dual_bound}
\proofsource{Analysis-of-Boolean-Functions-by-Ryan-ODonnell}{
  pdf-ffbe9c7b3306-p259-b002,
  pdf-ffbe9c7b3306-p259-b003
}
\label{FABL.selfAdjoint_dual_bound}
\leanok
\uses{FABL.exists_uniformLpNorm_dual, FABL.uniformInner, FABL.uniformInner_le_uniformLpNorm_mul, FABL.uniformLpNorm, FABL.uniformLpNorm_nonneg}
\difficulty{4}
Let $\Omega$ be a finite nonempty set, let $p,q,p',q',C\in\mathbb R$, and suppose $p'$
is Hölder-conjugate to $p$ and $q'$ is Hölder-conjugate to $q$. If a linear operator $T$
on real-valued functions on $\Omega$ is self-adjoint for the uniform inner product and
satisfies $\|Tf\|_q\le C\|f\|_p$ for every such function $f$, then $\|Tg\|_{p'}\le
C\|g\|_{q'}$ for every such function $g$, where these norms are taken with respect to
the uniform probability measure on $\Omega$.
\end{theorem}

\begin{definition}[Even coefficient of a sign function]
\lean{FABL.signEvenCoeff}
\statementsource{Analysis-of-Boolean-Functions-by-Ryan-ODonnell}{
  pdf-ffbe9c7b3306-p035-b008,
  pdf-ffbe9c7b3306-p035-b009,
  pdf-ffbe9c7b3306-p257-b004
}
\label{FABL.signEvenCoeff}
\leanok
\uses{FABL.Sign}
For a function $f:\{\pm1\}\to\mathbb{R}$, define its even coefficient to be
$\bigl(f(1)+f(-1)\bigr)/2$.
\end{definition}

\begin{definition}[Odd coefficient of a sign function]
\lean{FABL.signOddCoeff}
\statementsource{Analysis-of-Boolean-Functions-by-Ryan-ODonnell}{
  pdf-ffbe9c7b3306-p035-b008,
  pdf-ffbe9c7b3306-p035-b009,
  pdf-ffbe9c7b3306-p257-b004
}
\label{FABL.signOddCoeff}
\leanok
\uses{FABL.Sign}
For a function $f:\{\pm1\}\to\mathbb{R}$, define its odd coefficient to be
$\bigl(f(1)-f(-1)\bigr)/2$.
\end{definition}

\begin{theorem}[Affine reconstruction from sign values]
\lean{FABL.sign_affine_reconstruct}
\proofsource{Analysis-of-Boolean-Functions-by-Ryan-ODonnell}{pdf-ffbe9c7b3306-p035-b009}
\label{FABL.sign_affine_reconstruct}
\leanok
\uses{FABL.Sign, FABL.signEvenCoeff, FABL.signOddCoeff, FABL.signValue}
\difficulty{4}
Let $f$ be any real-valued function on the two signs $\{-1,1\}$, and let $x\in\{-1,1\}$.
Then $f(x)=\frac{f(1)+f(-1)}{2}+\frac{f(1)-f(-1)}{2}x$.
\end{theorem}

\begin{theorem}[Self-adjointness of the sign affine operator]
\lean{FABL.signAffine_selfAdjoint}
\label{FABL.signAffine_selfAdjoint}
\leanok
\uses{FABL.Sign, FABL.signEvenCoeff, FABL.signOddCoeff, FABL.signValue, FABL.uniformInner}
\difficulty{3}
For any real number $\rho$ and functions $f,g:\{-1,1\}\to\mathbb{R}$, the uniform inner
product of the affine transform of $f$ with $g$ equals the uniform inner product of $f$
with the affine transform of $g$. Here the transform sends $f(x)$ to
$\frac{f(1)+f(-1)}2+\rho\frac{f(1)-f(-1)}2x$.
\end{theorem}

\begin{theorem}[One-bit norm comparison]
\lean{FABL.oneBit_twoQ}
\proofsource{Analysis-of-Boolean-Functions-by-Ryan-ODonnell}{
  pdf-ffbe9c7b3306-p258-b001,
  pdf-ffbe9c7b3306-p258-b002
}
\label{FABL.oneBit_twoQ}
\leanok
\uses{FABL.Sign, FABL.exists_uniformLpNorm_dual, FABL.oneBit_pTwo, FABL.signAffine_selfAdjoint, FABL.signEvenCoeff, FABL.signOddCoeff, FABL.signValue, FABL.sign_affine_reconstruct, FABL.uniformInner, FABL.uniformInner_le_uniformLpNorm_mul, FABL.uniformLpNorm, FABL.uniformLpNorm_nonneg}
\difficulty{6}
Let $q,a,b,\rho\in\mathbb R$ with $2<q$, $0\le\rho$, and $\rho\le 1/\sqrt{q-1}$. For a
uniformly chosen sign $\varepsilon\in\{-1,1\}$, the $q$-norm of $a+\rho b\varepsilon$ is
at most the $2$-norm of $a+b\varepsilon$.
\end{theorem}

\begin{theorem}[Sharp one-bit norm comparison]
\lean{FABL.oneBit_twoQ_sharp}
\proofsource{Analysis-of-Boolean-Functions-by-Ryan-ODonnell}{pdf-ffbe9c7b3306-p259-b006}
\label{FABL.oneBit_twoQ_sharp}
\leanok
\uses{FABL.Sign, FABL.oneBit_twoQ, FABL.signValue, FABL.uniformLpNorm}
\difficulty{1}
Let $q,a,b\in\mathbb R$ with $2<q$. For a uniformly chosen sign
$\varepsilon\in\{-1,1\}$, the $q$-norm of $a+b\varepsilon/\sqrt{q-1}$ is at most the
$2$-norm of $a+b\varepsilon$.
\end{theorem}

\begin{theorem}[Two to q hypercontractivity]
\lean{FABL.uniformSign_twoQHypercontractive}
\proofsource{Analysis-of-Boolean-Functions-by-Ryan-ODonnell}{
  pdf-ffbe9c7b3306-p258-b001,
  pdf-ffbe9c7b3306-p258-b002,
  pdf-ffbe9c7b3306-p259-b006
}
\label{FABL.uniformSign_twoQHypercontractive}
\leanok
\uses{FABL.IsHypercontractiveRandomVariable, FABL.oneBit_twoQ_sharp, FABL.signValue, FABL.uniformLpNormENNReal, FABL.uniformLpNormENNReal_coe}
\difficulty{3}
Let $S=\{-1,1\}$ have the uniform probability distribution, let $X(s)=s$, and let $q$ be
a real number with $q>2$. For every $a,b\in\mathbb{R}$, the uniform $L^q$ norm of
$a+bX/\sqrt{q-1}$ is at most the uniform $L^2$ norm of $a+bX$.
\end{theorem}

\begin{theorem}[Dual hypercontractive bound]
\lean{FABL.selfAdjoint_fourThirdTwo_of_twoFour}
\proofsource{Analysis-of-Boolean-Functions-by-Ryan-ODonnell}{
  pdf-ffbe9c7b3306-p254-b003,
  pdf-ffbe9c7b3306-p254-b002
}
\label{FABL.selfAdjoint_fourThirdTwo_of_twoFour}
\leanok
\uses{FABL.selfAdjoint_dual_bound, FABL.uniformInner, FABL.uniformLpNorm}
\difficulty{3}
Let $n$ be a natural number and let $T$ be an operator on real-valued functions on the
sign cube $\{-1,1\}^n$. Suppose $T$ is self-adjoint for the uniform inner product and
satisfies $\|Tf\|_4\le\|f\|_2$ for every function $f$. Then, for every function $g$,
$\|Tg\|_2\le\|g\|_{4/3}$, with all norms taken under the uniform probability measure on
the cube.
\end{theorem}
